% GENERATED FILE. Edit canonical lessons, then run npm run generate:books.
% canonical-source-hash: fnv1a64:0f109ebbccdff464
% canonical-lessons: SA-C189-manjusa, SA-C189-sikyam, SA-C189-pitakam, SA-C189-katah, SA-C189-astaranam

\chapter{Trunk, Hanging basket, Basket, Mat, Carpet}
\label{ch:sa-trunk-hanging-basket-basket-mat-carpet}
\hlchaptermodality{189}

\begin{chapteropening}
\textbf{By the end of this chapter:} \emph{I can say a trunk, a hanging basket, a basket, a mat and a carpet.}

Complete the last lesson of chapter 189: I can say a trunk, a hanging basket, a basket, a mat and a carpet.
\end{chapteropening}

\section[\texorpdfstring{\sk{मञ्जूषा} (mañjūṣā)}{मञ्जूषा (mañjūṣā)}]{\sk{मञ्जूषा} (mañjūṣā) — a trunk, a chest}

\label{lesson:SA-C189-manjusa}



\begin{quote}
Before the new one: say the Sanskrit for spotless, then the Sanskrit for bright.
\end{quote}

\subsection*{You'll want to know: \sk{मञ्जूषा}}
\textbf{\sk{मञ्जूषा}} — \emph{mañjūṣā} — \textquotedblleft{}a trunk, a chest\textquotedblright{}.

\begin{grammarlens}[title={What you've built}]
One more word for this chapter.
\end{grammarlens}

\subsection*{Your turn}
\begin{itemize}
\raggedright
  \item \emph{Say it:} \emph{mañjūṣā}
  \item \emph{Say it:} \emph{mañjūṣā}, once more
  \item \emph{Say it:} say \emph{mañjūṣā}
  \item \emph{Recall:} say the Sanskrit for spotless, then the Sanskrit for bright, then say \emph{mañjūṣā} again
\end{itemize}

\begin{tcolorbox}[breakable,colback=teal!4,colframe=teal!35!black,title={Before you move on}]
What does \sk{मञ्जूषा} mean? (\textquotedblleft{}A trunk, a chest\textquotedblright{}.) Say it once more.
\end{tcolorbox}
\section[\texorpdfstring{\sk{शिक्यम्} (śikyam)}{शिक्यम् (śikyam)}]{\sk{शिक्यम्} (śikyam) — a hanging basket}

\label{lesson:SA-C189-sikyam}



\begin{quote}
Before the new one: say the Sanskrit for bright, then the Sanskrit for a trunk.
\end{quote}

\subsection*{You'll want to know: \sk{शिक्यम्}}
\textbf{\sk{शिक्यम्}} — \emph{śikyam} — \textquotedblleft{}a hanging basket\textquotedblright{}.

\begin{grammarlens}[title={What you've built}]
One more word for this chapter.
\end{grammarlens}

\subsection*{Your turn}
\begin{itemize}
\raggedright
  \item \emph{Say it:} \emph{śikyam}
  \item \emph{Say it:} \emph{śikyam}, once more
  \item \emph{Say it:} say \emph{śikyam}
  \item \emph{Recall:} say the Sanskrit for bright, then the Sanskrit for a trunk, then say \emph{śikyam} again
\end{itemize}

\begin{tcolorbox}[breakable,colback=teal!4,colframe=teal!35!black,title={Before you move on}]
What does \sk{शिक्यम्} mean? (\textquotedblleft{}A hanging basket\textquotedblright{}.) Say it once more.
\end{tcolorbox}
\section[\texorpdfstring{\sk{पिटकम्} (piṭakam)}{पिटकम् (piṭakam)}]{\sk{पिटकम्} (piṭakam) — a basket}

\label{lesson:SA-C189-pitakam}



\begin{quote}
Before the new one: say the Sanskrit for a trunk, then the Sanskrit for a hanging basket.
\end{quote}

\subsection*{You'll want to know: \sk{पिटकम्}}
\textbf{\sk{पिटकम्}} — \emph{piṭakam} — \textquotedblleft{}a basket\textquotedblright{}.

\begin{grammarlens}[title={What you've built}]
One more word for this chapter.
\end{grammarlens}

\subsection*{Your turn}
\begin{itemize}
\raggedright
  \item \emph{Say it:} \emph{piṭakam}
  \item \emph{Say it:} \emph{piṭakam}, once more
  \item \emph{Say it:} say \emph{piṭakam}
  \item \emph{Recall:} say the Sanskrit for a trunk, then the Sanskrit for a hanging basket, then say \emph{piṭakam} again
\end{itemize}

\begin{tcolorbox}[breakable,colback=teal!4,colframe=teal!35!black,title={Before you move on}]
What does \sk{पिटकम्} mean? (\textquotedblleft{}A basket\textquotedblright{}.) Say it once more.
\end{tcolorbox}
\section[\texorpdfstring{\sk{कटः} (kaṭaḥ)}{कटः (kaṭaḥ)}]{\sk{कटः} (kaṭaḥ) — a mat}

\label{lesson:SA-C189-katah}



\begin{quote}
Before the new one: say the Sanskrit for a hanging basket, then the Sanskrit for a basket.
\end{quote}

\subsection*{You'll want to know: \sk{कटः}}
\textbf{\sk{कटः}} — \emph{kaṭaḥ} — \textquotedblleft{}a mat\textquotedblright{}.

\begin{grammarlens}[title={What you've built}]
One more word for this chapter.
\end{grammarlens}

\subsection*{Your turn}
\begin{itemize}
\raggedright
  \item \emph{Say it:} \emph{kaṭaḥ}
  \item \emph{Say it:} \emph{kaṭaḥ}, once more
  \item \emph{Say it:} say \emph{kaṭaḥ}
  \item \emph{Recall:} say the Sanskrit for a hanging basket, then the Sanskrit for a basket, then say \emph{kaṭaḥ} again
\end{itemize}

\begin{tcolorbox}[breakable,colback=teal!4,colframe=teal!35!black,title={Before you move on}]
What does \sk{कटः} mean? (\textquotedblleft{}A mat\textquotedblright{}.) Say it once more.
\end{tcolorbox}
\section[\texorpdfstring{\sk{आस्तरणम्} (āstaraṇam)}{आस्तरणम् (āstaraṇam)}]{\sk{आस्तरणम्} (āstaraṇam) — a carpet, a spread}

\label{lesson:SA-C189-astaranam}



\begin{quote}
Before the new one: say the Sanskrit for a basket, then the Sanskrit for a mat.
\end{quote}

\subsection*{You'll want to know: \sk{आस्तरणम्}}
\textbf{\sk{आस्तरणम्}} — \emph{āstaraṇam} — \textquotedblleft{}a carpet, a spread\textquotedblright{}.

\begin{grammarlens}[title={What you've built}]
One more word for this chapter.
\end{grammarlens}

\subsection*{Your turn}
\begin{itemize}
\raggedright
  \item \emph{Say it:} \emph{āstaraṇam}
  \item \emph{Say it:} \emph{āstaraṇam}, once more
  \item \emph{Say it:} say \emph{āstaraṇam}
  \item \emph{Recall:} say the Sanskrit for a basket, then the Sanskrit for a mat, then say \emph{āstaraṇam} again
\end{itemize}

\begin{tcolorbox}[breakable,colback=teal!4,colframe=teal!35!black,title={Before you move on}]
What does \sk{आस्तरणम्} mean? (\textquotedblleft{}A carpet, a spread\textquotedblright{}.) Say it once more.
\end{tcolorbox}
